% Revize výukové přípravy, výchozí soubor GitHub Work-MAT, 95. hodina.
% Zdroje: Realisticky.cz 020203–020213, 020310–020320, 020409–020411; úlohy seřazené podle obtížnosti
% Úlohy jsou vlastní / úlohy výchozí přípravy, případně nově didakticky obměněné.
\documentclass[11pt,a4paper]{article}
\input{../../sablony/pripravy/priprava.sty}
\usepackage{tikz}
\nastavPripravu{1. ročník}{Rovnice a nerovnice}{95}{Souhrnné procvičování}
\begin{document}
\blok{Volitelné opakování před kvadratickými funkcemi}
Novou látku nezavádíme. Žáci si vybírají úlohy podle svých chyb z testů 86 a 94. K této přípravě patří zvláštní \textbf{sbírka úloh 95} s řešením pro učitele, sestavená na výběr, nikoli jako povinný test ze všech příkladů.
\blok{Organizace jedné vyučovací hodiny}
\begin{itemize}
\item 5 min: každý žák vybere podle předchozích chyb \emph{jednu slabší oblast}.
\item 7 min: společná oprava jedné typické chyby, např. změny znaménka při dělení.
\item 25 min: samostatná práce ze sbírky -- doporučené minimum: jedna rovnice, jedna nerovnice, jedna soustava a jedna absolutní hodnota. Silnější žáci přidají parametr a úlohu ze skupiny C.
\item 8 min: porovnání řešení, kontrola podmínek a pojmenování toho, co musí žák ještě procvičit.
\end{itemize}
\blok{Krátké řešené ukázky pro cílenou podporu}
\begin{prikladReseniBox}
Vyřešte $\frac{3x-2}{2}=\frac{x+4}{3}$.
\tcblower
\[\begin{aligned}
3(3x-2)&=2(x+4) \\
9x-6&=2x+8 \\
7x&=14 \\
x&=2 \\
K&=\{2\}
\end{aligned}\]

\end{prikladReseniBox}
\begin{prikladReseniBox}
Vyřešte $-3(x-1)>6$ a vysvětlete obrácení znaménka.
\tcblower
\[\begin{aligned}
-3x+3&>6 \\
-3x&>3 \\
x&<-1 \\
K&=(-\infty;-1)
\end{aligned}\]

\end{prikladReseniBox}
\blok{Příklady pro společnou kontrolu, dle potřeby}
\begin{zadaniReseniBox}
Vyřešte $\begin{cases}x+y=7,\\2x-y=2.\end{cases}$
\tcblower
\[\begin{aligned}
3x&=9 \\
x&=3,\quad y=4 \\
K&=\{[3;4]\}
\end{aligned}\]

\end{zadaniReseniBox}
\begin{zadaniReseniBox}
Vyřešte $|x-2|+|x+2|=6$.
\tcblower
\[\begin{aligned}
x<-2&:\ -2x=6\Rightarrow x=-3 \\
-2\le x\le2&:\ 4\ne6 \\
x>2&:\ 2x=6\Rightarrow x=3 \\
K&=\{-3;3\}
\end{aligned}\]

\end{zadaniReseniBox}
\begin{zadaniReseniBox}
Diskutujte $(a-2)x=2(a-2)$.
\tcblower
\[\begin{aligned}
a\ne2&:\ x=2 \\
a=2&:\ 0=0 \\
K_a&=\begin{cases}\{2\},&a\ne2,\\\mathbb R,&a=2.\end{cases}
\end{aligned}\]

\end{zadaniReseniBox}
\blok{Volitelná výzva}
\begin{zadaniReseniBox}
Kolik řešení má soustava $\begin{cases}2x+2y=8,\\x+y=4\end{cases}$?
\tcblower
\[\begin{aligned}
2(x+y)&=8\quad\text{-- stejné rovnice} \\
y&=4-x \\
K&=\{[t;4-t]\mid t\in\mathbb R\}
\end{aligned}\]
Graficky jde o dvě splývající přímky.
\end{zadaniReseniBox}
\blok{Domácí úkol}
\begin{zadaniReseniBox}[title={Domácí příprava}]
\begin{enumerate}[label=\alph*),itemsep=2pt]
\item $\dfrac{x+2}{3}=5$.
\item $|x-1|\le2$.
\item Z $P=2(a+b)$ vyjádřete $b$.
\item $\begin{cases}2x+y=7,\\x-y=2.\end{cases}$
\item $|x-3|+|x+1|=6$.
\item Diskutujte $(a+1)x=a+1$ pro všechna reálná $a$.
\end{enumerate}
\tcblower
\begin{enumerate}[label=\alph*),itemsep=1pt]
\item
\[\begin{aligned}
x+2&=15 \\
x&=13 \\
K&=\{13\}
\end{aligned}\]
\item
\[\begin{aligned}
-2&\le x-1\le2 \\
-1&\le x\le3 \\
K&=\langle-1;3\rangle
\end{aligned}\]
\item
\[\begin{aligned}
P&=2a+2b \\
2b&=P-2a \\
b&=\frac P2-a
\end{aligned}\]
\item
\[\begin{aligned}
3x&=9 \\
x&=3,\quad y=1 \\
K&=\{[3;1]\}
\end{aligned}\]
\item
\[\begin{aligned}
x<-1&:\ -2x+2=6\Rightarrow x=-2 \\
-1\le x\le3&:\ 4\ne6 \\
x>3&:\ 2x-2=6\Rightarrow x=4 \\
K&=\{-2;4\}
\end{aligned}\]
\item
\[\begin{aligned}
a\ne-1&:\ x=1 \\
a=-1&:\ 0=0 \\
K_a&=\begin{cases}\{1\},&a\ne-1,\\\mathbb R,&a=-1.\end{cases}
\end{aligned}\]
\end{enumerate}
\end{zadaniReseniBox}
\poznamka{Přílohu \texttt{95\_SB\_RovniceNerovniceSoustavy.pdf} žáci používají jako volitelnou sbírku. Nedávej všem všech 16 úloh: cílem je cílená samostatná práce a kontrola chyb. Následuje úvod do kvadratických funkcí.}
\end{document}
